% Part: normal-modal-logic
% Chapter: syntax-and-semantics
% Section: language-modal-logic

\documentclass[../../../include/open-logic-section]{subfiles}

\begin{document}

\olfileid[tr]{nml}{syn}{lan}

\olsection{Temel Kiplik Mantığının Dili}

\begin{defn}
Kiplik mantığının temel dili şunları içerir:
\begin{enumerate}
  \tagitem{prvFalse}{!!{falsity}~$\lfalse$ için önerme sabiti.}{}
  \tagitem{prvTrue}{!!{truth}~$\ltrue$ için önerme sabiti.}{}
  \item !!{denumerable}s bir !!{propositional variable}s kümesi: $\Obj
    p_0$, $\Obj p_1$, $\Obj p_2$, \dots
  \item Önerme bağlaçları: \startycommalist
  \iftag{prvNot}{\ycomma $\lnot$ (değilleme)}{}%
  \iftag{prvAnd}{\ycomma $\land$ (tümel evetleme)}{}%
  \iftag{prvOr}{\ycomma $\lor$ (tikel evetleme)}{}%
  \iftag{prvIf}{\ycomma $\lif$ (!!{conditional})}{}%
  \iftag{prvIff}{\ycomma $\liff$ (!!{biconditional})}{}.
  \tagitem{prvBox}{$\Box$ kiplik işleci.}{}
  \tagitem{prvDiamond}{$\Diamond$ kiplik işleci.}{}
\end{enumerate}
\end{defn}

\begin{defn}
Temel kiplik dilinin \emph{!!^{formula}s{}i} tümevarımlı olarak aşağıdaki gibi
tanımlanır:
\begin{enumerate}
\tagitem{prvFalse}{$\lfalse$ atomik bir !!{formula}{}dür.}{}

\tagitem{prvTrue}{$\ltrue$ atomik bir !!{formula}{}dür.}{}

\item Her !!{propositional variable} $\Obj p_i$, (atomik) bir !!{formula}{}dür.

\tagitem{prvNot}{$!A$ !!a{formula} ise $\lnot !A$ da
  !!a{formula}{}dür.}{}

\tagitem{prvAnd}{$!A$ ile $!B$ !!{formula}s ise $(!A \land
  !B)$ de !!a{formula}{}dür.}{}

\tagitem{prvOr}{$!A$ ile $!B$ !!{formula}s ise $(!A \lor !B)$
  de !!a{formula}{}dür.}{}

\tagitem{prvIf}{$!A$ ile $!B$ !!{formula}s ise $(!A \lif !B)$
  de !!a{formula}{}dür.}{}

\tagitem{prvIff}{$!A$ ile $!B$ !!{formula}s ise $(!A \liff !B)$
  de !!a{formula}{}dür.}{}

\tagitem{prvBox}{$!A$ !!a{formula} ise $\Box !A$ da
  !!a{formula}{}dür.}{}

\tagitem{prvDiamond}{$!A$ !!a{formula} ise $\Diamond !A$ da
  !!a{formula}{}dür.}{}

\tagitem{limitClause}{Başka hiçbir şey !!a{formula} değildir.}{}
\end{enumerate}
\end{defn}

\begin{tagblock}{defNot,defOr,defAnd,defIf,defIff,defTrue,defFalse,defBox,defDiamond}
\begin{defn}
Tanımlanmış işleçler kullanılarak kurulan formüller aşağıdaki gibi
anlaşılmalıdır:

\begin{tagenumerate}{defTrue,defFalse,defNot,defOr,defAnd,defIf,defIff,defBox,defDiamond}

\tagitem{defTrue}{$\ltrue$,
  \iftag{prvFalse}{$\lnot\lfalse$}{sabitlenmiş atomik bir
    !!{formula}~$!A$ için $(!A \lor \lnot !A)$} ifadesinin kısaltmasıdır.}{}

\tagitem{defFalse}{$\lfalse$,
  \iftag{prvTrue}{$\lnot\ltrue$}{sabitlenmiş atomik bir
    !!{formula}~$!A$ için $(!A \land \lnot !A)$} ifadesinin kısaltmasıdır.}{}

\tagitem{defNot}{$\lnot !A$, $!A \lif \lfalse$ ifadesinin kısaltmasıdır.}{}

\tagitem{defOr}{$!A \lor !B$,
  \iftag{prvAnd}{$\lnot(\lnot !A \land \lnot !B)$}{$\lnot !A \lif
    !B$} ifadesinin kısaltmasıdır.}{}

\tagitem{defAnd}{$!A \land !B$,
  \iftag{prvOr}{$\lnot(\lnot !A \lor \lnot !B)$}{$\lnot (!A \lif
    \lnot !B)$} ifadesinin kısaltmasıdır.}{}

\tagitem{defIf}{$!A \lif !B$,
  \iftag{prvOr}{$\lnot !A \lor !B$}{$\lnot (!A \land \lnot !B)$}
  ifadesinin kısaltmasıdır.}{}

\tagitem{defIff}{$!A \liff !B$, $(!A \lif !B) \land (!B
  \lif !A)$ ifadesinin kısaltmasıdır.}{}

\tagitem{defBox}{$\Box !A$, $\lnot\Diamond\lnot !A$ ifadesinin kısaltmasıdır. }{}

\tagitem{defDiamond}{$\Diamond !A$, $\lnot\Box\lnot !A$ ifadesinin kısaltmasıdır.}{}
\end{tagenumerate}
\end{defn}
\end{tagblock}

!!a{formula}~$!A$, $\Box$ ya da $\Diamond$ içermiyorsa ona
\emph{kipliksiz} deriz.

\end{document}
